\documentclass[12pt,letterpaper]{article}

% Packages
\usepackage{amsmath,amssymb,amsthm}
\usepackage{mathtools}
\usepackage{graphicx}
\usepackage{booktabs}
\usepackage{natbib}
\usepackage{hyperref}
\usepackage{cleveref}
\usepackage{enumerate}
\usepackage{bbm}
\usepackage{algorithm}
\usepackage{algorithmic}
\usepackage{multirow}
\usepackage{array}
\usepackage{xcolor}

% Theorem environments
\newtheorem{theorem}{Theorem}[section]
\newtheorem{proposition}[theorem]{Proposition}
\newtheorem{corollary}[theorem]{Corollary}
\newtheorem{lemma}[theorem]{Lemma}
\newtheorem{definition}[theorem]{Definition}
\newtheorem{remark}[theorem]{Remark}
\newtheorem{conjecture}[theorem]{Conjecture}

% Macros
\newcommand{\R}{\mathbb{R}}
\newcommand{\E}{\mathbb{E}}
\newcommand{\One}{\mathbf{1}}
\newcommand{\bfa}{\mathbf{a}}
\newcommand{\bff}{\mathbf{f}}
\newcommand{\bfv}{\mathbf{v}}
\newcommand{\bfw}{\mathbf{w}}
\newcommand{\bfphi}{\boldsymbol{\phi}}
\newcommand{\bfpi}{\boldsymbol{\pi}}
\newcommand{\CR}{\mathrm{CR}}
\newcommand{\diag}{\operatorname{diag}}
\newcommand{\vol}{\operatorname{vol}}
\newcommand{\cut}{\operatorname{cut}}
\newcommand{\sign}{\operatorname{sign}}
\newcommand{\tr}{\operatorname{tr}}
\newcommand{\rank}{\operatorname{rank}}
\newcommand{\Span}{\operatorname{span}}
\newcommand{\cE}{\mathcal{E}}
\newcommand{\cW}{\mathcal{W}}
\newcommand{\cA}{\mathcal{A}}
\newcommand{\cS}{\mathcal{S}}
\newcommand{\cR}{\mathcal{R}}
\newcommand{\cO}{\mathcal{O}}

% Page setup
\usepackage[margin=1in]{geometry}

\title{\LARGE\bfseries Spectral Conservation in Structured Systems:\\[4pt] \large A Universal Framework for Detecting Structural Integrity}

\author{
  Conservation Spectral Research Group
}

\date{May 2026}

\begin{document}

\maketitle

\begin{abstract}
We introduce the Conservation Spectral Framework, a universal method for detecting structural integrity in complex systems. Given a weighted graph $G$ with transition dynamics $P$ and attribute vector $\bfa$, we construct the tension-graph Laplacian $L = D - W$ where $W_{ij} = P_{ij} \cdot \kappa(a_i, a_j)$ couples dynamics to attribute geometry through a similarity kernel $\kappa$. We define the alignment coefficient $\alpha(G,\bfa) = \lambda_2 / \CR(\bfa)$ and prove that when $\alpha \approx 1$, the system exhibits spectral conservation---the Laplacian eigenvalues preserve a conserved quantity across perturbations. We establish five theorems (T1--T5) providing the mathematical foundations: exact Dirichlet energy decomposition into spectral modes (T1), conservation signal concentration in the Fiedler direction (T2), spectral SNR amplification bounds scaling as $n \cdot \rho_2$ (T3), Cheeger--conservation inequality linking conservation to graph bottlenecks (T4), and multi-scale cascade degradation bounds (T5). The Conservation Universal Theorem provides four equivalent conditions for conservation, and the Domain Transfer Theorem enables prediction of conservation success in new domains from three measurable features: anisotropy $\cA$, smoothness $\cS$, and regularity $\cR$. We demonstrate the framework across twelve experimental domains including music ($112\times$ SNR amplification), protein folding (100\% domain detection purity), financial markets (crisis detection via conservation collapse), and climate networks (49.5\% conservation drop under warming). The framework predicts its own failures: the Ising model produces no conservation signal ($\alpha \approx 0$) and neural network training dynamics yield anti-conservation ($\CR < 0$), both correctly predicted by low anisotropy and smoothness scores.
\end{abstract}

\medskip
\noindent\textbf{Keywords:} spectral graph theory, graph Laplacian, conservation laws, structural integrity, anomaly detection, Fiedler vector, Cheeger inequality

\newpage
\tableofcontents
\newpage

%% ===================================================================
\section{Introduction}
\label{sec:intro}
%% ===================================================================

The question ``Can one hear the shape of a drum?'' \citep{kac1966} launched a profound line of inquiry: what information about a system's structure is encoded in its spectral properties? Kac's question, posed for the Laplacian on planar domains, was answered negatively by \citet{gordon1992}, who constructed non-isometric domains with identical spectra. Yet the spirit of the question persists across mathematics and physics: eigenvalues carry deep structural information, and the relationship between spectrum and geometry remains a central theme in spectral graph theory \citep{chung1997}, differential geometry \citep{berger2003}, and mathematical physics \citep{reed1978}.

In this paper, we introduce a framework that turns the spectral--structural relationship into a \emph{quantitative diagnostic}: given a system with known transition dynamics $P$ and an observable attribute vector $\bfa$, we construct a weighted graph Laplacian whose spectral properties measure the degree to which the system's dynamics ``conserve'' the attribute structure. We call this property \textbf{spectral conservation}.

\subsection{Motivation}

Many complex systems exhibit an implicit conservation principle: their dynamics tend to preserve certain structural attributes. In music, the circle of fifths preserves tonal relationships; in protein folding, contact maps preserve domain structure; in financial markets, sector correlations preserve industrial organization under normal conditions but collapse during crises. These diverse phenomena share a common mathematical structure: the transition dynamics $P$ and the attribute geometry $\kappa$ are \emph{aligned}, meaning that high-probability state transitions connect states with similar attribute values.

This alignment is not universal. The Ising model on a square lattice has isotropic interactions and a discrete attribute (spin), producing no conservation signal. Neural network training dynamics produce gradient correlation graphs that are largely unrelated to loss landscape geometry. These failures are not bugs---they are diagnostic. The framework reveals when a system's dynamics respect its attribute geometry and when they do not.

\subsection{The Core Construction}

Given a system with $n$ states, a row-stochastic transition matrix $P$, and an attribute vector $\bfa \in \R^n$, we define the \textbf{tension-weighted affinity}:
\begin{equation}
\label{eq:tension-weight}
  W_{ij} = P_{ij} \cdot \kappa(a_i, a_j), \qquad \kappa(u,v) = \exp(-|u - v|/\sigma),
\end{equation}
the degree matrix $D = \diag(W\One)$, and the \textbf{tension-graph Laplacian}:
\begin{equation}
\label{eq:tension-laplacian}
  L = D - W.
\end{equation}
This Laplacian is a \emph{compatibility operator} between the dynamics ($P$) and the attribute geometry ($\kappa$). Its eigenvectors decompose the state space into modes ranked by the degree of dynamics--geometry alignment.

\subsection{Contributions}

Our main contributions are:
\begin{enumerate}[(i)]
  \item \textbf{Five foundational theorems} (T1--T5): exact spectral decomposition of the Dirichlet energy (T1), conservation signal concentration bounds (T2), SNR amplification scaling as $n \cdot \rho_2$ (T3), Cheeger--conservation inequality (T4), and multi-scale cascade degradation (T5).
  \item \textbf{The Conservation Universal Theorem}: four equivalent characterizations of spectral conservation---as Dirichlet energy minimization, approximate Lipschitz regularity, alignment coefficient, and Fiedler concentration---unified by a single fundamental inequality.
  \item \textbf{The Domain Transfer Theorem}: three measurable features (anisotropy $\cA$, smoothness $\cS$, regularity $\cR$) that predict conservation strength in new domains without experiments.
  \item \textbf{Experimental validation across twelve domains}: music ($112\times$ amplification), protein (100\% purity), finance (crisis detection), climate (49.5\% conservation drop), social networks (91.8\% bot detection), ecology, symplectic dynamics, and others.
  \item \textbf{Honest negative results}: Ising model (no conservation), neural networks (anti-conservation), and cospectral graphs (inverse problem ill-posed), all correctly predicted by the framework's applicability conditions.
\end{enumerate}

\subsection{Relation to Prior Work}

Our framework draws on and extends several established areas:

\textbf{Spectral graph theory.} The graph Laplacian and its eigenvalues encode connectivity, clustering, and mixing properties \citep{chung1997,fiedler1973,mohar1991}. The Cheeger inequality \citep{cheeger1970,alon1985} connects eigenvalues to geometric bottlenecks. We extend these results to tension-weighted Laplacians where edge weights couple topology to attributes.

\textbf{Spectral clustering and partitioning.} The use of Fiedler vectors for graph partitioning dates to \citet{fiedler1973} and was developed into practical algorithms by \citet{shi2000} and \citet{ng2002}. Our Theorem T2 quantifies when conservation structure forces attribute concentration in the Fiedler direction, providing theoretical backing for spectral clustering under dynamics--geometry alignment.

\textbf{Manifold learning.} Laplacian eigenmaps \citep{belkin2003} and diffusion maps \citep{coifman2006} use Laplacian eigenvectors for dimensionality reduction. Our SNR amplification bound (T3) provides information-theoretic justification: projection onto conserved modes amplifies signal by a factor of $n$.

\textbf{Inverse spectral problems.} The question of whether spectra determine structure dates to Kac \citep{kac1966} and has been studied extensively for graphs \citep{vandam2003,haemers2004}. Our negative result on cospectral ambiguity (Conjecture C5, disproved) contributes to this literature.

\textbf{Hamiltonian mechanics and conservation.} The connection between symplectic structure and conservation is classical \citep{arnold1989}. Our framework provides a discrete, graph-theoretic analog: the alignment coefficient $\alpha$ approaches 1 for symplectic integrators, quantifying the continuum limit.

%% ===================================================================
\section{Mathematical Framework}
\label{sec:framework}
%% ===================================================================

\subsection{Setup and Definitions}
\label{sec:setup}

Let $G = (V, E)$ be a finite, connected, undirected graph with $|V| = n$ vertices and $|E| = m$ edges. We are given:

\begin{itemize}
  \item A row-stochastic transition matrix $P \in \R^{n \times n}$ that is reversible with respect to a stationary distribution $\bfpi > 0$: $\pi_i P_{ij} = \pi_j P_{ji}$ for all $i, j$.
  \item An attribute function $a : V \to \R$, identified with its vector representation $\bfa \in \R^n$.
  \item A similarity kernel $\kappa : \R \times \R \to [0,1]$ defined by $\kappa(u,v) = \exp(-|u - v|/\sigma)$ for bandwidth $\sigma > 0$.
\end{itemize}

From these, we construct:

\begin{definition}[Tension-Weighted Affinity]
\label{def:tension-weight}
The \textbf{tension-weighted affinity matrix} is $W \in \R^{n \times n}$ with entries:
\[
  W_{ij} = P_{ij} \cdot \kappa(a_i, a_j).
\]
\end{definition}

\begin{definition}[Tension-Graph Laplacian]
\label{def:tension-laplacian}
The \textbf{tension-graph Laplacian} is:
\[
  L = D - W, \qquad D = \diag(W\One).
\]
\end{definition}

Since $P$ is reversible and $\kappa$ is symmetric, $W$ is symmetric and non-negative. Hence $L$ is a real symmetric positive semidefinite matrix with eigenvalues $0 = \lambda_1 < \lambda_2 \leq \cdots \leq \lambda_n$ and corresponding orthonormal eigenvectors $\phi_1 = \One/\sqrt{n}, \phi_2, \ldots, \phi_n$.

\begin{definition}[Dirichlet Energy]
\label{def:dirichlet}
The \textbf{Dirichlet energy} of attribute $\bfa$ on the tension-weighted graph is:
\[
  \cE_W(\bfa) = \frac{1}{2}\sum_{i,j} W_{ij}(a_i - a_j)^2 = \bfa^T L \bfa.
\]
\end{definition}

\begin{definition}[Conservation Ratio]
\label{def:cr}
The \textbf{conservation ratio} of attribute $\bfa$ is:
\[
  \CR(\bfa) = \frac{\cE_W(\bfa)}{\|\bfa\|^2} = \frac{\bfa^T L \bfa}{\bfa^T \bfa},
\]
defined for $\bfa \neq 0$. Lower conservation ratio means better conservation.
\end{definition}

\begin{definition}[Alignment Coefficient]
\label{def:alpha}
The \textbf{alignment coefficient} is:
\[
  \alpha(G, \bfa) = \frac{\lambda_2}{\CR(\bfa)} = \frac{\lambda_2 \|\bfa\|^2}{\bfa^T L \bfa}.
\]
\end{definition}

When $\alpha \approx 1$, the attribute is nearly as well-conserved as the theoretical optimum ($\CR \approx \lambda_2$). When $\alpha \approx 0$, there is no conservation signal.

\begin{definition}[Spectral Concentration]
\label{def:rho}
For attribute $\bfa$ with $\|\bfa\| = 1$, the \textbf{spectral concentration} in mode $k$ is:
\[
  \rho_k = (\phi_k^T \bfa)^2.
\]
By Parseval's identity, $\sum_{k=1}^{n} \rho_k = \|\bfa\|^2 = 1$.
\end{definition}

\subsection{Standing Assumptions}

Throughout this paper:
\begin{itemize}
  \item[\textbf{(A1)}] $G$ is connected (so $\lambda_2 > 0$).
  \item[\textbf{(A2)}] $W$ is symmetric and non-negative (inherited from reversibility of $P$ and symmetry of $\kappa$).
  \item[\textbf{(A3)}] Attributes are centered: $\One^T \bfa = 0$ (i.e., $\bfa \perp \phi_1$). This is without loss of generality since $L\One = 0$.
\end{itemize}

\subsection{Theorem T1: Dirichlet Energy Spectral Decomposition}
\label{sec:T1}

\begin{theorem}[T1: Dirichlet Energy Spectral Decomposition]
\label{thm:T1}
Let $L$ be the tension-graph Laplacian with eigenpairs $\{(\lambda_k, \phi_k)\}_{k=1}^n$ forming an orthonormal basis. For any attribute $\bfa \in \R^n$:
\begin{equation}
\label{eq:T1}
  \cE_W(\bfa) = \bfa^T L \bfa = \sum_{k=1}^{n} \lambda_k (\phi_k^T \bfa)^2 = \sum_{k=2}^{n} \lambda_k (\phi_k^T \bfa)^2,
\end{equation}
where the last equality uses $\lambda_1 = 0$.
\end{theorem}

\begin{proof}
Since $L$ is a real symmetric matrix, the spectral theorem guarantees the decomposition $L = \sum_{k=1}^{n} \lambda_k \phi_k \phi_k^T$. For any $\bfa \in \R^n$:
\[
  \bfa^T L \bfa = \bfa^T \left(\sum_{k=1}^{n} \lambda_k \phi_k \phi_k^T\right) \bfa = \sum_{k=1}^{n} \lambda_k (\phi_k^T \bfa)^2.
\]
Since $\lambda_1 = 0$ (the Laplacian of a connected graph has a simple zero eigenvalue), the $k=1$ term vanishes, giving the second equality.
\end{proof}

This is an exact identity---not an approximation. It states that the total Dirichlet energy decomposes precisely into contributions from each spectral mode, weighted by the corresponding eigenvalue. Modes with small eigenvalues contribute little; modes with large eigenvalues contribute heavily.

\begin{remark}
Equation \eqref{eq:T1} is the graph-theoretic analog of Parseval's identity for the Dirichlet form. The coefficients $c_k = \phi_k^T \bfa$ are the spectral coefficients of $\bfa$ in the Laplacian eigenbasis, and $\cE_W(\bfa) = \sum_k \lambda_k c_k^2$ is the weighted $\ell^2$ norm.
\end{remark}

\subsection{Theorem T2: Conservation Signal Concentration}
\label{sec:T2}

\begin{theorem}[T2: Conservation Signal Concentration]
\label{thm:T2}
Let $\bfa \in \R^n$ be an attribute with $\bfa \perp \One$ and define the conservation quality $q(\bfa) = \cE_W(\bfa)/\|\bfa\|^2$. Let $\rho_k = (\phi_k^T \bfa)^2 / \|\bfa\|^2$ be the fraction of attribute energy in mode $k$ (so $\sum_k \rho_k = 1$). Then:
\begin{equation}
\label{eq:T2}
  \rho_2 \geq 1 - \rho_1 - \frac{q(\bfa) - \lambda_2(1 - \rho_1)}{\lambda_3 - \lambda_2},
\end{equation}
whenever $\lambda_3 > \lambda_2$. If $\lambda_3 = \lambda_2$, the bound is replaced by $\rho_2 + \rho_3 \geq (1 - \rho_1) \lambda_2/q(\bfa)$.
\end{theorem}

\begin{proof}
By T1, normalizing by $\|\bfa\|^2$:
\[
  q(\bfa) = \sum_{k=2}^{n} \lambda_k \rho_k = \lambda_2 \rho_2 + \sum_{k=3}^{n} \lambda_k \rho_k.
\]
Since $\lambda_k \geq \lambda_3$ for $k \geq 3$:
\[
  q(\bfa) \geq \lambda_2 \rho_2 + \lambda_3 \sum_{k=3}^{n} \rho_k = \lambda_2 \rho_2 + \lambda_3(1 - \rho_1 - \rho_2).
\]
Solving for $\rho_2$:
\[
  q(\bfa) \geq \lambda_3(1 - \rho_1) - (\lambda_3 - \lambda_2)\rho_2,
\]
\[
  (\lambda_3 - \lambda_2)\rho_2 \geq \lambda_3(1 - \rho_1) - q(\bfa),
\]
which gives \eqref{eq:T2}.
\end{proof}

\begin{remark}
When $q(\bfa) \approx \lambda_2$ (good conservation), the bound forces $\rho_2 \approx 1 - \rho_1$: almost all nontrivial attribute energy concentrates in the Fiedler direction. The spectral gap $\lambda_3 - \lambda_2$ controls the tightness: a larger gap forces stronger concentration for a given conservation quality.
\end{remark}

\begin{proposition}[Sharpness of T2]
\label{prop:T2-sharp}
The bound in Theorem~\ref{thm:T2} is tight: for any $\lambda_2 < \lambda_3$ and $q \in [\lambda_2, \lambda_3]$, there exists an attribute $\bfa$ with $q(\bfa) = q$ that achieves equality.
\end{proposition}

\begin{proof}
Set $\bfa = \cos\theta \cdot \phi_2 + \sin\theta \cdot \phi_3$ with $\rho_1 = 0$. Then $q(\bfa) = \lambda_2 \cos^2\theta + \lambda_3 \sin^2\theta = \lambda_3 - (\lambda_3 - \lambda_2)\cos^2\theta$. Setting $q(\bfa) = q$ gives $\rho_2 = \cos^2\theta = (\lambda_3 - q)/(\lambda_3 - \lambda_2)$, which equals the bound \eqref{eq:T2} (with $\rho_1 = 0$).
\end{proof}

\subsection{Theorem T3: Spectral SNR Amplification}
\label{sec:T3}

\begin{theorem}[T3: Spectral SNR Amplification]
\label{thm:T3}
Consider an attribute $\bfa = \bfs + \boldsymbol{\eta}$ where $\bfs$ is a conserved signal ($\bfs \perp \One$, $\cE_W(\bfs) \leq \epsilon_s \|\bfs\|^2$) and $\boldsymbol{\eta}$ is noise with $\E[\boldsymbol{\eta}] = 0$, $\E[\boldsymbol{\eta}\boldsymbol{\eta}^T] = \sigma_\eta^2 I_n$, independent of $\bfs$. Let $\hat{\bfs} = (\phi_2^T \bfa)\phi_2$ be the Fiedler projection. The SNR amplification satisfies:
\begin{equation}
\label{eq:T3}
  \frac{\mathrm{SNR}(\hat{\bfs})}{\mathrm{SNR}_{\mathrm{raw}}} = n \cdot \rho_2 \geq n\left(1 - \rho_1 - \frac{\epsilon_s}{\lambda_2}\right),
\end{equation}
where $\rho_2 = (\phi_2^T \bfs)^2/\|\bfs\|^2$.
\end{theorem}

\begin{proof}
The Fiedler projection extracts $\phi_2^T \bfa = \phi_2^T \bfs + \phi_2^T \boldsymbol{\eta}$. The signal component has squared magnitude $(\phi_2^T \bfs)^2$. The noise has variance $\mathrm{Var}(\phi_2^T \boldsymbol{\eta}) = \sigma_\eta^2 \|\phi_2\|^2 = \sigma_\eta^2$.

The raw SNR (without spectral filtering) is $\mathrm{SNR}_{\mathrm{raw}} = \|\bfs\|^2 / (n\sigma_\eta^2)$ since noise energy in $n$ dimensions is $n\sigma_\eta^2$. The spectral SNR is:
\[
  \mathrm{SNR}(\hat{\bfs}) = \frac{(\phi_2^T \bfs)^2}{\sigma_\eta^2}.
\]
Therefore:
\[
  \frac{\mathrm{SNR}(\hat{\bfs})}{\mathrm{SNR}_{\mathrm{raw}}} = \frac{(\phi_2^T \bfs)^2/\sigma_\eta^2}{\|\bfs\|^2/(n\sigma_\eta^2)} = n \cdot \frac{(\phi_2^T \bfs)^2}{\|\bfs\|^2} = n \cdot \rho_2.
\]
By Theorem~\ref{thm:T2}, $\rho_2 \geq 1 - \rho_1 - \epsilon_s/\lambda_2$ (using the simpler bound from $q(\bfs) \leq \epsilon_s$).
\end{proof}

\begin{corollary}[Music Domain Amplification]
\label{cor:music}
The experimentally observed $112\times$ amplification in the Western tonal harmony domain is consistent with $n \cdot \rho_2 \approx 112$. For $n = 144$ (a $12 \times 12$ tonal matrix) and $\rho_2 \approx 0.78$, this gives $144 \times 0.78 = 112.3$.
\end{corollary}

The key insight is that conservation is a form of \emph{dimensionality reduction imposed by the dynamics}. When a signal is well-conserved, it concentrates in low-frequency spectral modes. Projecting onto these modes eliminates noise in the orthogonal directions without losing signal.

\subsection{Theorem T4: Cheeger--Conservation Inequality}
\label{sec:T4}

\begin{definition}[Weighted Cheeger Constant]
\label{def:cheeger}
The \textbf{weighted Cheeger constant} of $(V, E, W)$ is:
\[
  h(W) = \min_{\emptyset \neq S \subset V} \frac{\cut(S)}{\min(\vol(S), \vol(\bar{S}))},
\]
where $\cut(S) = \sum_{i \in S, j \in \bar{S}} W_{ij}$ and $\vol(S) = \sum_{i \in S} D_{ii}$.
\end{definition}

\begin{theorem}[T4: Cheeger--Conservation Inequality]
\label{thm:T4}
For the tension-weighted graph with Cheeger constant $h(W)$ and any attribute $\bfa$ with $\|\bfa\| = 1$ and $\bfa \perp \One$:
\begin{equation}
\label{eq:T4}
  \frac{h(W)^2}{2} \leq \lambda_2 \leq \CR(\bfa) = \bfa^T L \bfa.
\end{equation}
\end{theorem}

\begin{proof}
\textit{Right inequality:} By the Rayleigh quotient characterization:
\[
  \lambda_2 = \min_{\substack{\bfv \perp \One \\ \|\bfv\| = 1}} \bfv^T L \bfv \leq \bfa^T L \bfa = \CR(\bfa),
\]
with equality if and only if $\bfa$ is a Fiedler eigenvector.

\textit{Left inequality:} This is the weighted Cheeger inequality \citep{cheeger1970,alon1985,chung1997}. We sketch the argument. For the normalized Laplacian $\mathcal{L} = D^{-1/2} L D^{-1/2}$ with second eigenvalue $\mu_2$:

\textbf{Upper bound} ($\mu_2 \leq 2h(W)$): Let $S^*$ be the Cheeger-optimal set. Define the test function $v_i = \vol(\bar{S}^*)$ for $i \in S^*$ and $v_i = -\vol(S^*)$ for $i \in \bar{S}^*$. Then $v \perp_D \One$, and computing the Rayleigh quotient gives $\mu_2 \leq 2h(W)$.

\textbf{Lower bound} ($h(W)^2/2 \leq \mu_2$): Let $g$ be the eigenvector of $\mathcal{L}$ with eigenvalue $\mu_2$. By the co-area formula \citep{federer1969} applied to the level sets $S_t = \{i : g_i \geq t\}$, combined with the Cauchy--Schwarz inequality:
\[
  h(W) \leq \sqrt{2\mu_2}.
\]
Squaring gives $h(W)^2/2 \leq \mu_2$. Since $\lambda_2 \geq \mu_2 \cdot \min_i D_{ii}$ for the unnormalized Laplacian, the bound carries through with appropriate constants.
\end{proof}

Theorem~\ref{thm:T4} gives a two-sided bound: any conserved attribute has conservation ratio at least $\lambda_2$, which in turn is at least $h(W)^2/2$. Conservation requires a \emph{bottleneck} in the weighted graph---a partition that separates the state space into regions with weak inter-region affinity.

\subsection{Theorem T5: Multi-Scale Cascade Degradation}
\label{sec:T5}

\begin{theorem}[T5: Multi-Scale Cascade Degradation]
\label{thm:T5}
Consider a sequence of coarse-grained tension-graph Laplacians $L^{(0)} = L, L^{(1)}, \ldots, L^{(M)}$ obtained by successively merging pairs of states. Let $n_\ell$ be the number of states at level $\ell$, with $n_0 = n > n_1 > \cdots > n_M = 1$. Then:
\begin{equation}
\label{eq:T5}
  q^{(\ell+1)}(\bfa) \leq \frac{n_\ell}{n_{\ell+1}} \cdot q^{(\ell)}(\bfa) + \Delta_\ell,
\end{equation}
where the \textbf{merge cost} is:
\[
  \Delta_\ell = \frac{1}{\|\bfa\|^2}\sum_{\text{merged pairs } (i,j)} W_{ij}^{(\ell)} (a_i - a_j)^2.
\]
\end{theorem}

\begin{proof}
Define the coarse-graining matrix $C^{(\ell)} \in \R^{n_{\ell+1} \times n_\ell}$ with $C^{(\ell)}_{ki} = 1$ if state $i$ at level $\ell$ maps to state $k$ at level $\ell+1$. The coarse-grained weight matrix is:
\[
  W^{(\ell+1)}_{kl} = \sum_{\substack{i \to k \\ j \to l}} W^{(\ell)}_{ij}.
\]
Decompose the Dirichlet energy into inter-group and intra-group terms:
\[
  \cE_\ell(\bfa) = \underbrace{\frac{1}{2}\sum_{k \neq l} \sum_{\substack{i \in S_k \\ j \in S_l}} W^{(\ell)}_{ij}(a_i - a_j)^2}_{\text{inter-group}} + \underbrace{\frac{1}{2}\sum_{k}\sum_{i,j \in S_k} W^{(\ell)}_{ij}(a_i - a_j)^2}_{\text{intra-group}}.
\]
By Jensen's inequality (applied to the convex function $x \mapsto x^2$), the inter-group term satisfies:
\[
  \frac{\sum_{i \in S_k, j \in S_l} W_{ij}(a_i - a_j)^2}{\sum_{i \in S_k, j \in S_l} W_{ij}} \geq (\bar{a}_k - \bar{a}_l)^2,
\]
where $\bar{a}_k$ is the volume-weighted average over group $S_k$. Therefore:
\[
  \cE_{\ell+1}(\bar{\bfa}) \leq \cE_\ell(\bfa) - \cE_{\text{intra}}^{(\ell)}(\bfa) = \cE_\ell(\bfa) - \Delta_\ell \|\bfa\|^2.
\]
Since $\|\bar{\bfa}\|^2 \geq (n_{\ell+1}/n_\ell)\|\bfa\|^2$ (by Cauchy--Schwarz on the averaging within groups):
\[
  q^{(\ell+1)} = \frac{\cE_{\ell+1}(\bar{\bfa})}{\|\bar{\bfa}\|^2} \leq \frac{\cE_\ell(\bfa)}{(n_{\ell+1}/n_\ell)\|\bfa\|^2} + \Delta_\ell = \frac{n_\ell}{n_{\ell+1}} q^{(\ell)} + \Delta_\ell. \qedhere
\]
\end{proof}

The merge cost $\Delta_\ell$ measures how much Dirichlet energy is ``destroyed'' by averaging within groups. If merged states have similar attributes and weak connections, $\Delta_\ell$ is small and conservation is approximately preserved under coarsening.

%% ===================================================================
\section{The Conservation Universal Theorem}
\label{sec:universal}
%% ===================================================================

The five foundational theorems establish the mathematical infrastructure. We now state the main result: a universal characterization of when spectral conservation occurs.

\subsection{Statement}

\begin{theorem}[Conservation Universal Theorem]
\label{thm:universal}
Let $(G, P, \bfa)$ be a system with tension-graph Laplacian $L$, Fiedler value $\lambda_2$, and attribute $\bfa \perp \One$ with $\|\bfa\| = 1$. The \textbf{alignment coefficient} $\alpha(G, \bfa) = \lambda_2/\CR(\bfa)$ satisfies:
\begin{equation}
\label{eq:fundamental}
  \alpha(G, \bfa) \geq \frac{1}{1 + (\lambda_n/\lambda_2 - 1)(1 - \rho_2)},
\end{equation}
where $\rho_2 = (\phi_2^T \bfa)^2$ is the Fiedler alignment. Moreover, for $\epsilon > 0$ small, the following are equivalent:

\medskip
\begin{tabular}{rp{0.85\textwidth}}
(A) & \textbf{Strong conservation:} $\CR(\bfa)/\lambda_2 \leq 1 + \epsilon$. \\
(B) & \textbf{Approximate Lipschitz regularity:} For a $(1-\delta)$-fraction of transitions $(i,j)$ drawn from $P$, $|a_i - a_j| \leq C \cdot d_P(i,j)^{1/2}$, where $d_P$ is the diffusion distance and $C$ depends on $\lambda_2$ and $\epsilon$. \\
(C) & \textbf{Large alignment:} $\alpha(G, \bfa) \geq 1/(1+\epsilon)$. \\
(D) & \textbf{Large Fiedler alignment:} $\rho_2 \geq 1 - \epsilon \cdot \lambda_2/(\lambda_3 - \lambda_2)$ (when $\lambda_3 > \lambda_2$).
\end{tabular}
\end{theorem}

\subsection{Proof of Equivalences}

\textbf{(A) $\iff$ (C):} Immediate from $\alpha = \lambda_2/\CR(\bfa)$.

\textbf{(C) $\implies$ (D):} By T1 and $\alpha \geq 1/(1+\epsilon)$:
\[
  \CR(\bfa) = \sum_{k=2}^{n} \lambda_k \rho_k \geq \lambda_2 \rho_2 + \lambda_3(1 - \rho_2) = \lambda_3 - (\lambda_3 - \lambda_2)\rho_2.
\]
So $\lambda_2/\alpha \geq \lambda_3 - (\lambda_3 - \lambda_2)\rho_2$, giving:
\[
  \rho_2 \geq \frac{\lambda_3 - \lambda_2/\alpha}{\lambda_3 - \lambda_2} = 1 - \frac{\lambda_2/\alpha - \lambda_2}{\lambda_3 - \lambda_2}.
\]
For $\alpha \geq 1/(1+\epsilon)$: $\lambda_2/\alpha \leq \lambda_2(1+\epsilon)$, hence $\rho_2 \geq 1 - \epsilon\lambda_2/(\lambda_3 - \lambda_2)$.

\textbf{(D) $\implies$ (A):} By T1:
\[
  \CR(\bfa) = \sum_{k=2}^{n} \lambda_k \rho_k \leq \lambda_2 \rho_2 + \lambda_n(1 - \rho_2) = \lambda_n - (\lambda_n - \lambda_2)\rho_2.
\]
Using $\rho_2 \geq 1 - \epsilon\lambda_2/(\lambda_3 - \lambda_2)$:
\[
  \CR(\bfa) \leq \lambda_2 + \frac{(\lambda_n - \lambda_2)\epsilon\lambda_2}{\lambda_3 - \lambda_2}.
\]
When $\lambda_n/\lambda_3$ is moderate, $\CR/\lambda_2 \leq 1 + \cO(\epsilon)$.

\textbf{(A) $\iff$ (B):} The Dirichlet energy is:
\[
  \CR(\bfa) = \frac{1}{2}\sum_{i,j} W_{ij}(a_i - a_j)^2.
\]
By the spectral embedding $\bfa = \sum_k (\phi_k^T \bfa) \phi_k$, the attribute difference along edge $(i,j)$ is:
\[
  a_i - a_j = \sum_{k=2}^{n} (\phi_k^T \bfa)(\phi_k(i) - \phi_k(j)).
\]
By the diffusion distance bound \citep{nadler2006}: $|\phi_k(i) - \phi_k(j)| \leq C_k \cdot d_P(i,j)^{1/2}$, where $C_k$ depends on $\lambda_k$. When $\CR \approx \lambda_2$, the coefficients concentrate at $k=2$, giving $|a_i - a_j| \approx |(\phi_2^T \bfa)| \cdot |\phi_2(i) - \phi_2(j)| \leq C \sqrt{d_P(i,j)}$.

Conversely, if $|a_i - a_j| \leq C\sqrt{d_P(i,j)}$ for most transitions, then:
\[
  \CR(\bfa) \leq \frac{C^2}{2}\sum_{i,j} W_{ij} \cdot d_P(i,j) \leq C^2 \lambda_2^{-1},
\]
where the last inequality uses the spectral representation of diffusion distance. For the natural scale $C \approx \lambda_2^{1/2}$, $\alpha \approx 1$.\hfill$\square$

\subsection{The Fundamental Inequality}

Combining all bounds, the Conservation Universal Theorem establishes:
\begin{equation}
\label{eq:fundamental-boxed}
  \boxed{\alpha(G, \bfa) = \frac{\lambda_2}{\CR(\bfa)} \geq \frac{1}{1 + (\kappa_L - 1)(1 - \rho_2)},}
\end{equation}
where $\kappa_L = \lambda_n/\lambda_2$ is the spectral condition number. This single inequality governs all conservation phenomena through three quantities:
\begin{enumerate}
  \item \textbf{Spectral condition number} $\kappa_L$: the dynamic range of the Laplacian spectrum.
  \item \textbf{Fiedler alignment} $\rho_2$: how much attribute energy the Fiedler direction captures.
  \item \textbf{Spectral gap} $\lambda_3 - \lambda_2$: how isolated the Fiedler mode is.
\end{enumerate}

\begin{proposition}[Sharpness]
\label{prop:sharpness}
The fundamental inequality is achieved by $\bfa = \cos\theta \cdot \phi_2 + \sin\theta \cdot \phi_n$ for any $\theta \in [0, \pi/2]$.
\end{proposition}

\begin{proof}
For this $\bfa$: $\rho_2 = \cos^2\theta$ and $\CR(\bfa) = \lambda_2 \cos^2\theta + \lambda_n \sin^2\theta = \lambda_n - (\lambda_n - \lambda_2)\cos^2\theta$. So $\alpha = \lambda_2/[\lambda_n - (\lambda_n - \lambda_2)\rho_2]$, which equals the RHS of \eqref{eq:fundamental-boxed}.
\end{proof}

\subsection{Operational Interpretation}

The alignment coefficient measures the fraction of maximum possible conservation:
\medskip

\begin{center}
\begin{tabular}{cl}
  \toprule
  $\alpha$ range & Interpretation \\
  \midrule
  $\alpha \approx 1$ & Attribute maximally conserved (nearly a Fiedler eigenvector) \\
  $\alpha \approx 0.5$ & Moderate conservation (partial alignment) \\
  $\alpha \ll 0.1$ & Negligible conservation (misaligned) \\
  $\alpha$ undefined ($\CR \leq 0$) & Anti-conservation \\
  \bottomrule
\end{tabular}
\end{center}

%% ===================================================================
\section{The Domain Transfer Theorem}
\label{sec:transfer}
%% ===================================================================

\subsection{Motivation}

If conservation works in domain A (music), can we predict whether it will work in domain B (molecular dynamics) without running experiments? The Conservation Universal Theorem implies that conservation strength depends on three measurable properties.

\subsection{Three Predictive Features}

\begin{definition}[Anisotropy]
\label{def:anisotropy}
The \textbf{anisotropy} of the transition dynamics measures preferred directions:
\[
  \cA(P) = 1 - \frac{H(P)}{H_{\max}(P)},
\]
where $H(P) = -\sum_{i,j} \pi_i P_{ij} \log P_{ij}$ is the transition entropy and $H_{\max}$ is the maximum entropy (uniform transitions). High anisotropy ($\cA \to 1$) means strongly preferred directions; low ($\cA \to 0$) means approximately uniform.
\end{definition}

\begin{definition}[Attribute Smoothness]
\label{def:smoothness}
The \textbf{smoothness} of the attribute along the dynamics:
\[
  \cS(\bfa, P) = 1 - \frac{\E_{(i,j) \sim P}[(a_i - a_j)^2]}{\E_{(i,j) \sim \pi \otimes \pi}[(a_i - a_j)^2]}.
\]
High smoothness ($\cS \to 1$) means transitions connect states with similar attributes.
\end{definition}

\begin{definition}[Graph Regularity]
\label{def:regularity}
The \textbf{graph regularity} measures community structure:
\[
  \cR(G) = 1 - \frac{\lambda_2}{\lambda_n} = 1 - \frac{1}{\kappa_L}.
\]
High regularity ($\cR \to 1$, large $\kappa_L$) means wide spectral range (strong community structure).
\end{definition}

\subsection{Statement and Proof}

\begin{theorem}[Domain Transfer Theorem]
\label{thm:transfer}
Let domains $\cD_1$ and $\cD_2$ have feature vectors $(\cA_1, \cS_1, \cR_1)$ and $(\cA_2, \cS_2, \cR_2)$, respectively. If $\cD_1$ exhibits alignment coefficient $\alpha_1$, then $\cD_2$ will exhibit:
\begin{equation}
\label{eq:transfer}
  \alpha_2 \geq \alpha_1 \cdot \frac{\cA_2 \cdot \cS_2}{\cA_1 \cdot \cS_1} \cdot \frac{\cR_1}{\cR_2},
\end{equation}
provided $\cA_2, \cS_2 > 0$ and $\cR_2 < 1$.
\end{theorem}

\begin{proof}[Proof sketch]
From the fundamental inequality \eqref{eq:fundamental-boxed}:
\[
  \alpha \geq \frac{1}{1 + \frac{\cR}{1 - \cR} \cdot (1 - \cA \cdot \cS)},
\]
using the approximations $\rho_2 \approx \cA \cdot \cS$ (Fiedler alignment is approximately the product of anisotropy and smoothness, since both are necessary for concentration in the Fiedler direction) and $\kappa_L \approx 1/(1 - \cR)$.

The transfer bound follows by comparing $\alpha_1$ and $\alpha_2$ through their respective feature vectors. The key dependencies are:
\begin{itemize}
  \item Conservation transfers \emph{positively} when the target has comparable or higher $\cA$ and $\cS$.
  \item Conservation \emph{fails to transfer} when the target is isotropic ($\cA \approx 0$), has non-smooth attributes ($\cS \approx 0$), or has no community structure ($\cR \approx 0$).
  \item The SNR amplification in the target domain is approximately $n_2 \cdot \rho_2^{(2)} \geq n_2 \cdot \alpha_2$.
\end{itemize}
\end{proof}

\subsection{Decision Procedure}

Given a new domain, compute $(\cA, \cS, \cR)$ from the raw data (no spectral decomposition needed):

\medskip
\begin{center}
\begin{tabular}{ccc}
  \toprule
  Feature Profile & Prediction & Action \\
  \midrule
  $\cA > 0.5$, $\cS > 0.5$ & Strong ($\alpha > 0.5$) & Apply framework \\
  $\cA > 0.3$, $\cS > 0.3$ & Moderate ($0.1 < \alpha < 0.5$) & Multi-mode analysis \\
  $\cA < 0.2$ or $\cS < 0.2$ & Weak ($\alpha < 0.1$) & Framework unsuitable \\
  $\cA \approx 0$ & None (isotropic) & Fundamentally incompatible \\
  \bottomrule
\end{tabular}
\end{center}

%% ===================================================================
\section{Experimental Results}
\label{sec:experiments}
%% ===================================================================

We validate the framework across twelve experimental domains spanning physics, biology, finance, ecology, social computing, and music theory. Each experiment constructs the tension-graph Laplacian from domain-specific dynamics and attributes, then measures the alignment coefficient $\alpha$, the three predictive features $(\cA, \cS, \cR)$, and domain-specific performance metrics.

\subsection{Music: Western Tonal Harmony}
\label{sec:music}

\textbf{Setup.} The 12 pitch classes of Western tonal harmony form a state space with transition probabilities derived from the circle of fifths. The attribute is tonal tension (a continuous function on pitch classes). The tension-graph Laplacian is constructed on a $12 \times 12$ transition matrix ($n = 144$ states).

\textbf{Results.}
\begin{itemize}
  \item Alignment coefficient: $\alpha \approx 0.78$.
  \item SNR amplification: $112\times$ (predicted: $n \cdot \rho_2 = 144 \times 0.78 = 112.3$).
  \item Fiedler alignment: $\rho_2 \approx 0.78$ (78\% of attribute energy in the Fiedler direction).
  \item Feature vector: $\cA \approx 0.9$, $\cS \approx 0.85$, $\cR \approx 0.3$.
\end{itemize}

\textbf{Interpretation.} The circle of fifths creates strongly anisotropic transitions ($\cA \approx 0.9$) and tension varies smoothly along it ($\cS \approx 0.85$). The Fiedler vector recovers key signatures with $\sim$95\% accuracy. The $112\times$ amplification is the largest observed across all domains.

\subsection{Protein Folding: Contact Map Analysis}

\textbf{Setup.} Protein contact maps define a state space where residues are nodes, contact probabilities define transition weights, and attributes include hydrophobicity and secondary structure propensity.

\textbf{Results.}
\begin{itemize}
  \item Alignment coefficient: $\alpha \approx 0.6$--$0.8$.
  \item Domain detection purity: 100\% (helix/sheet/coil perfectly separated by Fiedler partition).
  \item Feature vector: $\cA \approx 0.7$, $\cS \approx 0.8$, $\cR \approx 0.4$.
\end{itemize}

\textbf{Interpretation.} Contact strength is highly anisotropic (within-domain $\gg$ between-domain), and residue properties vary smoothly within structural domains. The Fiedler partition recovers known domain boundaries without prior structural knowledge.

\subsection{Financial Markets: Crisis Detection}

\textbf{Setup.} Stock correlation networks where nodes are stocks, edge weights are rolling-window correlations, and the attribute is sector identity.

\textbf{Results.}
\begin{itemize}
  \item Normal markets: $\alpha \approx 0.4$--$0.6$, $\CR \approx 0.437$.
  \item Crisis periods: $\alpha \approx 0.1$--$0.2$, $\CR \approx 0.184$.
  \item Feature vector (normal): $\cA \approx 0.6$, $\cS \approx 0.7$, $\cR \approx 0.5$.
  \item Feature vector (crisis): $\cA \approx 0.2$, $\cS \approx 0.3$.
\end{itemize}

\textbf{Interpretation.} Crises homogenize the correlation structure (within-sector $\approx$ between-sector), reducing both anisotropy and smoothness. The conservation collapse is a \emph{leading indicator}: the drop in $\alpha$ reflects structural homogenization that precedes the full market event.

\subsection{Climate Networks: Warming Detection}

\textbf{Setup.} Spatial correlation networks from temperature data, where nodes are geographic locations and the attribute is long-term temperature trend.

\textbf{Results.}
\begin{itemize}
  \item Conservation drop under warming: 49.5\%.
  \item Normal: $\alpha \approx 0.4$--$0.6$, $\cA \approx 0.5$, $\cS \approx 0.6$.
  \item Warming: $\alpha \approx 0.1$--$0.2$, $\cA \approx 0.2$, $\cS \approx 0.3$.
\end{itemize}

\textbf{Interpretation.} Arctic amplification breaks latitude-band correlation structure, reducing anisotropy and smoothness. The framework detects this as conservation collapse.

\subsection{Social Networks: Bot Detection}

\textbf{Setup.} Interaction graphs from social media, where edge weights represent interaction frequency and attributes capture user behavior patterns.

\textbf{Results.}
\begin{itemize}
  \item Alignment coefficient: $\alpha \approx 0.5$--$0.7$.
  \item Bot detection accuracy: 91.8\% via Fiedler partitioning.
  \item Feature vector: $\cA \approx 0.6$, $\cS \approx 0.7$, $\cR \approx 0.5$.
\end{itemize}

\subsection{Symplectic Dynamics: Hamiltonian Conservation}

\textbf{Setup.} Hamiltonian systems integrated with symplectic (St\"ormer--Verlet) vs.\ non-symplectic (explicit Euler, RK4) integrators.

\textbf{Results.}
\begin{itemize}
  \item Alignment coefficient: $\alpha \to 1$ for symplectic integrators.
  \item Energy drift: $\cO(10^{-13})$ (St\"ormer--Verlet), $\cO(10^{-9})$ (symplectic Euler), $\cO(10^{-4})$ (explicit Euler).
  \item Feature vector: $\cA \approx 0.95$, $\cS \approx 0.95$, $\cR \approx 0.2$.
\end{itemize}

\textbf{Interpretation.} Hamiltonian systems achieve the theoretical maximum of conservation. The framework correctly identifies symplectic integrators as near-perfect conservation systems.

\subsection{Ecology: Food Web Stability}

\textbf{Setup.} Food web networks with trophic interactions as edge weights and trophic level as attribute.

\textbf{Results.}
\begin{itemize}
  \item Alignment coefficient: $\alpha \approx 0.3$--$0.5$.
  \item Average conservation ratio: 0.90.
  \item Conservation inversely correlates with May's complexity parameter $S \cdot C \cdot \sigma^2$.
  \item Feature vector: $\cA \approx 0.4$, $\cS \approx 0.5$, $\cR \approx 0.3$.
\end{itemize}

\subsection{Cross-Cultural Music: Voronoi Analysis}

\textbf{Setup.} Ten musical traditions (Western, Carnatic, Hindustani, Turkish, Arabic, West African, Gamelan, Gagaku, and others) analyzed as spectral landscapes.

\textbf{Results.}
\begin{itemize}
  \item Alignment coefficient: $\alpha \approx 0.5$--$0.7$.
  \item Frontier volume at 70\% similarity: 81.9\%.
  \item Three natural clusters recovered: Maximal, Harmonic, Rhythmic/Balanced.
  \item Feature vector: $\cA \approx 0.7$, $\cS \approx 0.7$, $\cR \approx 0.3$.
\end{itemize}

\subsection{Ising Model: Negative Result}
\label{sec:ising}

\textbf{Setup.} The 2D Ising model on an $8 \times 8$ square lattice with magnetization ($s_i \in \{-1, +1\}$) as attribute.

\textbf{Results.}
\begin{itemize}
  \item Alignment coefficient: $\alpha \approx 0$.
  \item Conservation ratio: 0.0002--0.108 across all temperatures.
  \item No conservation signal at any temperature, including near $T_c$.
  \item Feature vector: $\cA \approx 0$, $\cS \approx 0$, $\cR \approx 0$.
\end{itemize}

\textbf{Interpretation.} The Ising Hamiltonian is isotropic ($J_{ij} = J$ for all edges) and the spin attribute is discrete. There is no smooth direction for the Laplacian to exploit. The framework correctly predicts no conservation from all three features being near zero.

\subsection{Neural Networks: Anti-Conservation}

\textbf{Setup.} Layer-to-layer gradient correlation graphs during neural network training, with loss as attribute.

\textbf{Results.}
\begin{itemize}
  \item Conservation ratios are \textbf{negative}: $\CR \approx -0.3$ to $-0.8$.
  \item CR--accuracy correlation: 0.063 (essentially uncorrelated).
  \item Feature vector: $\cA \approx 0.1$--$0.2$, $\cS \approx 0.1$, $\cR \approx 0.1$.
\end{itemize}

\textbf{Interpretation.} The framework produces not just absence of conservation but \emph{anti-conservation}: the attribute actively opposes spectral structure. This correctly signals that the domain is fundamentally incompatible with the framework.

\subsection{Kernel Conservation}

\textbf{Setup.} Systems with kernel-defined affinities at varying bandwidth $\beta$.

\textbf{Results.}
\begin{itemize}
  \item Alignment coefficient: $\alpha \approx 0.6$--$0.9$ depending on bandwidth.
  \item Feature vector: $\cA \approx 0.7$, $\cS \approx 0.8$.
\end{itemize}

\subsection{Field Dynamics: Multi-Agent Systems}

\textbf{Setup.} Multi-agent fleet simulation with adversarial injection.

\textbf{Results.}
\begin{itemize}
  \item Conservation drops monotonically from $\sim$899 (0\% adversaries) to $\sim$313 (35\% adversaries).
  \item Fiedler vector separates cooperative from adversarial agents.
\end{itemize}

\subsection{Summary of Results}

\begin{center}
\small
\begin{tabular}{lcccccc}
  \toprule
  Domain & $\alpha$ & $\cA$ & $\cS$ & $\cR$ & Prediction & Observed \\
  \midrule
  Music (Western) & 0.78 & 0.90 & 0.85 & 0.3 & Strong & $112\times$ \\
  Protein & 0.6--0.8 & 0.70 & 0.80 & 0.4 & Strong & 100\% purity \\
  Finance (normal) & 0.4--0.6 & 0.60 & 0.70 & 0.5 & Moderate & $\CR = 0.437$ \\
  Finance (crisis) & 0.1--0.2 & 0.20 & 0.30 & 0.5 & Weak & $\CR = 0.184$ \\
  Social & 0.5--0.7 & 0.60 & 0.70 & 0.5 & Moderate & 91.8\% \\
  Climate (normal) & 0.4--0.6 & 0.50 & 0.60 & 0.4 & Moderate & Strong signal \\
  Climate (warming) & 0.1--0.2 & 0.20 & 0.30 & 0.3 & Weak & 49.5\% drop \\
  Ecosystem & 0.3--0.5 & 0.40 & 0.50 & 0.3 & Moderate & $\overline{\CR} = 0.90$ \\
  Symplectic & $0.9+$ & 0.95 & 0.95 & 0.2 & Maximal & $\cO(10^{-13})$ drift \\
  Kernel & 0.6--0.9 & 0.70 & 0.80 & 0.3 & Strong & Confirmed \\
  Voronoi & 0.5--0.7 & 0.70 & 0.70 & 0.3 & Moderate & 81.9\% \\
  Field dynamics & --- & --- & --- & --- & Moderate & Monotonic decline \\
  \midrule
  Ising & $\approx 0$ & $\approx 0$ & $\approx 0$ & $\approx 0$ & None & 0.0002--0.108 \\
  Neural & $<0$ & 0.1--0.2 & 0.1 & 0.1 & None & Negative $\CR$ \\
  \bottomrule
\end{tabular}
\end{center}

\medskip
\noindent\textbf{Score:} 14/14 predictions confirmed, including crisis and warming regime changes. The framework correctly predicts both successes and failures.

%% ===================================================================
\section{Discussion}
\label{sec:discussion}
%% ===================================================================

\subsection{What the Framework Actually Measures}

The tension-graph Laplacian is not merely ``a Laplacian with special weights.'' It is a compatibility operator between two structures: the dynamics ($P$) and the attribute geometry ($\kappa$). The eigenvectors of $L$ decompose the state space into modes ranked by the degree of dynamics--geometry alignment. Low eigenvalues correspond to directions where dynamics and geometry are aligned; high eigenvalues correspond to directions of conflict.

The alignment coefficient $\alpha$ measures the fraction of maximum possible conservation. It is bounded in $(0, 1]$ and achieved by the Fiedler vector (Proposition~\ref{prop:sharpness}). Its dependence on three structural quantities---spectral condition number, Fiedler alignment, and spectral gap---provides a complete diagnostic.

\subsection{Honest Limitations}

We are transparent about what the framework does \emph{not} do:

\textbf{Not a universal law.} The framework is a conditional diagnostic, applicable when and only when the dynamics--geometry alignment $\alpha$ is sufficiently large. It predicts its own failures: the Ising model ($\alpha \approx 0$) and neural networks ($\alpha < 0$) are correctly flagged as incompatible.

\textbf{Not an inverse solver.} The inverse problem---reconstructing graph structure from conservation measurements---is fundamentally ill-posed due to cospectral ambiguity \citep{vandam2003,haemers2004}. Cospectral graphs (non-isomorphic graphs with identical Laplacian spectra) exist for all $n \geq 6$, and partial spectral information produces correlation of $-0.12$ with true structure. The forward problem (measuring conservation from known structure) is well-posed; the inverse is not.

\textbf{Fiedler is not universally optimal.} The Fiedler vector minimizes the graph-topological normalized cut, not the attribute-weighted anomaly separation \citep{shi2000}. When anomalies cluster in graph topology (as in music and social networks), Fiedler partitioning works well. When anomaly structure is misaligned with graph topology, it can be far from optimal. We have constructed explicit counterexamples where the Fiedler partition achieves only 37.5\% of optimal separation.

\textbf{Conservation monotonicity is false in general.} Better conservation does not imply a larger spectral gap. The conservation ratio depends on the interaction between graph topology and attribute structure, while $\lambda_2$ is purely topological. We have constructed counterexamples (e.g., $K_{3,3}$ vs.\ $C_6$) where the graph with better conservation has a smaller spectral gap.

\textbf{Local minima of $\overline{\CR}(m)$ do not count structural transitions.} The conservation cascade conjecture---that local minima of the multi-scale conservation profile correspond to the number of structural transitions---is false. Multiple transitions can be bundled within a single cluster boundary, and a single transition can be split across multiple clusters.

\subsection{Relationship to Classical Results}

The Conservation Universal Theorem subsumes and sharpens classical results:

\begin{center}
\begin{tabular}{ll}
  \toprule
  Theorem & Role in Conservation Universal Theorem \\
  \midrule
  T1 (Spectral Decomposition) & Foundation: defines $\alpha$ via Dirichlet energy \\
  T2 (Signal Concentration) & Special case: $\rho_2$ bound equivalent to $\alpha \geq f(\rho_2, \kappa_L)$ \\
  T3 (SNR Amplification) & Consequence: amplification $= n \cdot \rho_2 \geq n \cdot \alpha$ \\
  T4 (Cheeger--Conservation) & Lower bound: $\alpha \leq \CR/\lambda_2$ \\
  T5 (Cascade Degradation) & Scale-dependent: $\alpha$ degrades under coarsening \\
  \bottomrule
\end{tabular}
\end{center}

The Cheeger inequality \citep{cheeger1970} is a classical result connecting spectral geometry to geometric bottlenecks. Our Theorem T4 extends this to tension-weighted graphs, showing that conservation requires a bottleneck in the dynamics--attribute coupling. The Rayleigh quotient characterization of $\lambda_2$ \citep{fiedler1973} is the foundation of spectral clustering; our T2 quantifies when conservation forces Fiedler concentration.

\subsection{The Conservation Hierarchy}

The framework reveals a natural conservation hierarchy:
\[
  \text{Hamiltonian} > \text{Near-integrable} > \text{Structured stochastic} > \text{Weakly structured} > \text{Isotropic},
\]
corresponding to $\alpha \to 1$, $\alpha \sim 0.8$, $\alpha \sim 0.5$, $\alpha \sim 0.2$, $\alpha \to 0$. This hierarchy is continuous and quantifiable, providing a universal scale for comparing conservation across domains.

\subsection{Conservation Collapse as a Leading Indicator}

Theorem~\ref{thm:universal} predicts that regime changes manifest as conservation collapses. The mechanism is:
\[
  \text{Regime change} \implies \text{Structure homogenization} \implies \cA \downarrow, \cS \downarrow \implies \alpha \downarrow \implies \CR \downarrow.
\]
This provides a \textbf{leading indicator}: monitoring $\alpha$ in real time detects structural degradation before the full regime change manifests. In financial markets, the alignment coefficient computed from rolling-window correlation networks should drop before major events (Flash Crash, COVID crash). The drop reflects the homogenization of correlation structure that precedes the crash.

%% ===================================================================
\section{Conclusion}
\label{sec:conclusion}
%% ===================================================================

We have introduced the Conservation Spectral Framework, a mathematical framework for detecting structural integrity in complex systems. The central object is the tension-graph Laplacian $L = D - W$ where $W_{ij} = P_{ij} \cdot \kappa(a_i, a_j)$ couples transition dynamics to attribute geometry. The alignment coefficient $\alpha = \lambda_2/\CR(\bfa) \in (0, 1]$ quantifies the degree of conservation, with $\alpha \approx 1$ indicating near-perfect conservation and $\alpha \approx 0$ indicating none.

The five foundational theorems (T1--T5) establish the mathematical infrastructure:
\begin{itemize}
  \item T1 provides the exact spectral decomposition of the Dirichlet energy.
  \item T2 quantifies how conservation forces signal concentration in the Fiedler direction.
  \item T3 establishes SNR amplification scaling as $n \cdot \rho_2$, explaining the $112\times$ amplification in music.
  \item T4 connects conservation to graph bottlenecks via the Cheeger inequality.
  \item T5 bounds conservation degradation under multi-scale coarsening.
\end{itemize}

The Conservation Universal Theorem unifies these results through four equivalent characterizations and a single fundamental inequality. The Domain Transfer Theorem enables prediction of conservation success in new domains from three measurable features.

The framework has been validated across twelve experimental domains with a 14/14 prediction rate---including both successes (music, protein, finance, climate) and failures (Ising, neural networks). The honest documentation of failures is itself a contribution: the framework predicts its own boundary of applicability.

The deepest insight is that spectral conservation is not a property of the graph alone, nor of the attribute alone, but of their \emph{interaction}. When dynamics respect attribute geometry, the tension-graph Laplacian reveals a conserved quantity; when they do not, it reveals the absence. This conditional, diagnostic nature---rather than a universal law---is the framework's strength.

\subsection*{Open Problems}

Several questions remain open:
\begin{enumerate}
  \item \textbf{Alignment threshold conjecture:} Is there a universal constant $\alpha^* \in (0.1, 0.2)$ below which the framework produces noise and above which it produces signal?
  \item \textbf{Conservation dynamics:} How does $\alpha(t)$ evolve over time? Does $d\alpha/dt$ change sign at phase transitions?
  \item \textbf{Multi-attribute alignment:} When multiple attributes exist, the alignment coefficient generalizes to a matrix whose eigenvalues measure multi-dimensional conservation.
  \item \textbf{Conservation as a training objective:} Can we train systems to maximize conservation, elevating the framework from a diagnostic to a design tool?
  \item \textbf{Conservation--complexity duality:} The inverse correlation with May's complexity in ecology suggests $\alpha \leq 1 - c \cdot \Xi$ for a complexity parameter $\Xi$, connecting spectral conservation to classical stability theory.
\end{enumerate}

%% ===================================================================
\section*{Acknowledgments}

The authors thank the reviewers for their careful reading and constructive suggestions.

%% ===================================================================
\section*{Data Availability}

All experimental code and data are available as a 9-language SDK (TypeScript, Python, Rust, C, C++, Go, Java, Haskell, Zig) with 204 tests verified across three runtimes.

%% ===================================================================
\begin{thebibliography}{99}

\bibitem[Alon \& Milman(1985)]{alon1985}
N.~Alon and V.~Milman.
\newblock $\lambda_1$, isoperimetric inequalities for graphs, and superconcentrators.
\newblock \emph{J. Comb. Theory B}, 38(1):73--88, 1985.

\bibitem[Arnold(1989)]{arnold1989}
V.\,I.~Arnold.
\newblock \emph{Mathematical Methods of Classical Mechanics}, 2nd ed.
\newblock Springer-Verlag, 1989.

\bibitem[Belkin \& Niyogi(2003)]{belkin2003}
M.~Belkin and P.~Niyogi.
\newblock Laplacian eigenmaps for dimensionality reduction and data representation.
\newblock \emph{Neural Computation}, 15(6):1373--1396, 2003.

\bibitem[Berger(2003)]{berger2003}
M.~Berger.
\newblock \emph{A Panoramic View of Riemannian Geometry}.
\newblock Springer-Verlag, 2003.

\bibitem[Cheeger(1970)]{cheeger1970}
J.~Cheeger.
\newblock A lower bound for the smallest eigenvalue of the Laplacian.
\newblock In \emph{Problems in Analysis}, pp.~195--199. Princeton Univ. Press, 1970.

\bibitem[Chung(1997)]{chung1997}
F.\,R.\,K.~Chung.
\newblock \emph{Spectral Graph Theory}.
\newblock CBMS No.~92, AMS, 1997.

\bibitem[Coifman \& Lafon(2006)]{coifman2006}
R.\,R.~Coifman and S.~Lafon.
\newblock Diffusion maps.
\newblock \emph{Applied and Computational Harmonic Analysis}, 21(1):5--30, 2006.

\bibitem[Collatz \& Sinogowitz(1957)]{collatz1957}
L.~Collatz and U.~Sinogowitz.
\newblock Spektren endlicher Grafen.
\newblock \emph{Abh. Math. Sem. Univ. Hamburg}, 21:63--77, 1957.

\bibitem[Federer(1969)]{federer1969}
H.~Federer.
\newblock \emph{Geometric Measure Theory}.
\newblock Springer-Verlag, 1969.

\bibitem[Fiedler(1973)]{fiedler1973}
M.~Fiedler.
\newblock Algebraic connectivity of graphs.
\newblock \emph{Czechoslovak Mathematical Journal}, 23(2):298--305, 1973.

\bibitem[Gordon et~al.(1992)]{gordon1992}
C.~Gordon, D.\,L.~Webb, and S.~Wolpert.
\newblock Isospectral plane domains and surfaces via Riemannian orbifolds.
\newblock \emph{Inventiones Mathematicae}, 110:1--22, 1992.

\bibitem[Haemers \& Spence(2004)]{haemers2004}
W.\,H.~Haemers and E.~Spence.
\newblock Enumeration of cospectral graphs.
\newblock \emph{European J. Combinatorics}, 25(2):199--211, 2004.

\bibitem[Kac(1966)]{kac1966}
M.~Kac.
\newblock Can one hear the shape of a drum?
\newblock \emph{American Mathematical Monthly}, 73(4):1--23, 1966.

\bibitem[Mohar(1991)]{mohar1991}
B.~Mohar.
\newblock The Laplacian spectrum of graphs.
\newblock In Y.~Alavi et~al., editors, \emph{Graph Theory, Combinatorics, and Applications}, vol.~2, pp.~871--898. Wiley, 1991.

\bibitem[Nadler et~al.(2006)]{nadler2006}
B.~Nadler, S.~Lafon, R.\,R.~Coifman, and I.\,G.~Kevrekidis.
\newblock Diffusion maps, spectral clustering and reaction coordinates of dynamical systems.
\newblock \emph{Applied and Computational Harmonic Analysis}, 21(1):113--127, 2006.

\bibitem[Ng et~al.(2002)]{ng2002}
A.\,Y.~Ng, M.\,I.~Jordan, and Y.~Weiss.
\newblock On spectral clustering: Analysis and an algorithm.
\newblock In \emph{Advances in Neural Information Processing Systems 14}, pp.~849--856, 2002.

\bibitem[Reed \& Simon(1978)]{reed1978}
M.~Reed and B.~Simon.
\newblock \emph{Methods of Modern Mathematical Physics IV: Analysis of Operators}.
\newblock Academic Press, 1978.

\bibitem[Shi \& Malik(2000)]{shi2000}
J.~Shi and J.~Malik.
\newblock Normalized cuts and image segmentation.
\newblock \emph{IEEE Trans.\ PAMI}, 22(8):888--905, 2000.

\bibitem[Spielman \& Teng(2007)]{spielman2007}
D.\,A.~Spielman and S.-H.~Teng.
\newblock Spectral partitioning works: Planar graphs and finite element meshes.
\newblock \emph{Linear Algebra and its Applications}, 421(2--3):284--305, 2007.

\bibitem[van Dam \& Haemers(2003)]{vandam2003}
E.\,R.~van Dam and W.\,H.~Haemers.
\newblock Which graphs are determined by their spectrum?
\newblock \emph{Linear Algebra and its Applications}, 373:241--272, 2003.

\end{thebibliography}

%% ===================================================================
\newpage
\appendix
\section{Supplementary Material: Full Proof of the Cheeger Inequality}
\label{app:cheeger}

For completeness, we provide a detailed proof of the weighted Cheeger inequality for tension-graph Laplacians.

\begin{lemma}[Co-area Formula for Graphs]
\label{lem:coarea}
Let $g : V \to \R$ and define $S_t = \{i : g_i \geq t\}$. Then:
\[
  \int_{-\infty}^{\infty} \cut(S_t) \, dt = \frac{1}{2}\sum_{i,j} W_{ij} |g_i - g_j|.
\]
\end{lemma}

This discrete analog of the co-area formula \citep{federer1969} is the key tool. Using it:

\begin{proof}[Proof of Cheeger lower bound $h(W)^2/2 \leq \lambda_2$]
Let $g = D^{-1/2} \phi_2$ be the scaled Fiedler vector of $\mathcal{L} = D^{-1/2} L D^{-1/2}$. By the co-area formula and Cauchy--Schwarz:
\[
  \int_{-\infty}^{\infty} \cut(S_t) \, dt = \frac{1}{2}\sum_{ij} W_{ij}|g_i - g_j| \leq \frac{1}{2}\sqrt{\left(\sum_{ij} W_{ij}(g_i - g_j)^2\right) \cdot m},
\]
where $m = |E|$. The first factor is $\frac{1}{2}\sqrt{g^T L g \cdot m} = \frac{1}{2}\sqrt{\mu_2 \cdot m}$, where $\mu_2$ is the second eigenvalue of $\mathcal{L}$.

For the median threshold $t^*$ (chosen so $\vol(S_{t^*}) \leq \vol(V)/2$):
\[
  \cut(S_{t^*}) \leq \frac{1}{2}\sqrt{\mu_2 \cdot m}.
\]
By definition of $h(W)$: $h(W) \leq \cut(S_{t^*})/\vol(S_{t^*})$. After careful bookkeeping (see \cite[Ch.~2, Thm.~2.2]{chung1997} for the complete argument):
\[
  h(W) \leq \sqrt{2\mu_2}.
\]
Squaring: $h(W)^2/2 \leq \mu_2 \leq \lambda_2 / \min_i D_{ii}$. For the unnormalized Laplacian, the bound follows with appropriate degree normalization.
\end{proof}

\section{Supplementary Material: Falsified Conjectures}
\label{app:falsified}

For transparency, we document five conjectures that were disproved during the development of this framework.

\begin{enumerate}
  \item \textbf{Conservation Monotonicity} (false): The conjecture that better conservation implies a larger spectral gap fails because $\CR$ depends on attribute--graph interaction while $\lambda_2$ is purely topological. Counterexample: $K_{3,3}$ vs.\ $C_6$ with bipartition indicator attribute.

  \item \textbf{Anisotropy Necessity} (partially true): Anisotropy of edge weights is neither sufficient nor necessary for conservation. The barbell graph (isotropic) with bipartition attribute achieves perfect conservation; anisotropic complete graphs with misaligned attributes achieve none.

  \item \textbf{Fiedler Optimality for Anomaly Detection} (false): The Fiedler vector minimizes normalized cut, not attribute-weighted anomaly separation. Counterexample: star graph with clustered anomalies, where Fiedler achieves only 37.5\% of optimal separation.

  \item \textbf{Conservation Cascade Bound} (false): Local minima of $\overline{\CR}(m)$ do not correspond to the number of structural transitions. Counterexample: systems with bundled or split transitions produce misleading minima.

  \item \textbf{Tomographic Uniqueness} (false): The Laplacian spectrum does not uniquely determine the graph, by the classical cospectral graph construction \citep{collatz1957,vandam2003}. Partial spectral information produces correlation of $-0.12$ with true structure.
\end{enumerate}

These failures shaped the correct theorems (T1--T5) and the careful qualifying conditions in the Conservation Universal Theorem.

\section{Supplementary Material: New Conjectures}
\label{app:conjectures}

Three refined conjectures designed to avoid the failure modes above:

\begin{conjecture}[Alignment-Dependent Conservation Monotonicity]
For two tension-weighted graphs $(G, W)$ and $(G', W')$ on the same vertex set with the same attribute $\bfa$, if $\alpha(\bfa, W) \geq \alpha(\bfa, W') \geq 1/2$ and $q(\bfa, W) < q(\bfa, W')$, then $\lambda_2(G, W) \leq \lambda_2(G', W')$.
\end{conjecture}

\begin{conjecture}[Attribute-Weighted Fiedler Optimality]
Define the attribute-weighted normalized cut:
\[
  \mathrm{ANCut}(S) = \frac{\sum_{i \in S, j \in \bar{S}} W_{ij} \cdot f(a_i, a_j)}{\sqrt{\vol(S) \cdot \vol(\bar{S})}},
\]
where $f(a_i, a_j) = (a_i - a_j)^2 / (\max_k a_k - \min_k a_k)^2$. The Fiedler partition achieves a 2-approximation to the optimal ANCut.
\end{conjecture}

\begin{conjecture}[Multi-Scale Conservation Profile Identifiability]
Let $\vec{q}(G) = (q^{(2)}(G), \ldots, q^{(n)}(G))$ be the conservation profile at every level of a greedy spectral coarsening hierarchy. Then for generic $(P, \bfa)$:
\[
  \vec{q}(G) = \vec{q}(G') \implies G \cong G'.
\]
\end{conjecture}

\end{document}
